\documentclass[border=8pt]{standalone}
\usepackage{fontspec}
\setmainfont{Noto Sans}[Ligatures=TeXOff]
\XeTeXgenerateactualtext=1
\setsansfont{Noto Sans}
\newfontfamily\bengalifont{Noto Sans Bengali}[Script=Bengali, NFSSFamily=bnfam]
\usepackage[Latin,Bengali]{ucharclasses}
% group-free font switching (as in header-common.tex): robust inside braces, cells and quotes
\makeatletter
\def\BN@saved{\rmdefault}
\def\BN@bn{bnfam}
\newcommand\BNin{\edef\BN@tmp{\f@family}\ifx\BN@tmp\BN@bn\else\let\BN@saved\BN@tmp\fi\fontfamily{bnfam}\selectfont}
\newcommand\BNout{\fontfamily{\BN@saved}\selectfont}
\makeatother
\setTransitionsFor{Bengali}{\BNin}{\BNout}
% Typographic quotes and dashes are not in a ucharclasses class (class 0), so a Bangla word right
% after “ or ‘ would not trigger the Bangla font. Treat them as Latin so that every switch fires.
\XeTeXcharclass"2018=\XeTeXcharclass`A
\XeTeXcharclass"2019=\XeTeXcharclass`A
\XeTeXcharclass"201C=\XeTeXcharclass`A
\XeTeXcharclass"201D=\XeTeXcharclass`A
\XeTeXcharclass"2013=\XeTeXcharclass`A
\XeTeXcharclass"2014=\XeTeXcharclass`A
\XeTeXcharclass"2026=\XeTeXcharclass`A
\usepackage{newunicodechar}
\newunicodechar{।}{{\bengalifont\char"0964}}
\usepackage{tikz}
\usetikzlibrary{arrows.meta,positioning,shapes.geometric,shapes.misc,shapes.symbols,fit,backgrounds,calc,chains}
\definecolor{ink}{HTML}{111827}
\definecolor{band}{HTML}{F3F4F6}
\definecolor{band2}{HTML}{E5E7EB}
\definecolor{accent}{HTML}{1D4ED8}
\definecolor{accentlight}{HTML}{DBEAFE}
\definecolor{warn}{HTML}{B91C1C}
\definecolor{warnlight}{HTML}{FEE2E2}
\definecolor{ok}{HTML}{047857}
\definecolor{oklight}{HTML}{D1FAE5}
\tikzset{
  >={Stealth[length=2.6mm]},
  box/.style={draw=ink, line width=0.6pt, rounded corners=9pt, fill=white, align=center,
              inner sep=4pt, font=\small},
  pill/.style={box, rounded corners=13pt},
  key/.style={box, fill=accentlight, draw=accent, line width=0.9pt},
  soft/.style={box, fill=band},
  dec/.style={draw=ink, line width=0.6pt, diamond, aspect=2.2, align=center, inner sep=1.5pt,
              font=\footnotesize, fill=white},
  arr/.style={->, draw=ink, line width=0.6pt},
  darr/.style={->, draw=ink, line width=0.6pt, dashed},
  lbl/.style={font=\footnotesize, fill=white, inner sep=1.5pt},
  layer/.style={draw=band2, fill=band, rounded corners=6pt},
  title/.style={font=\bfseries\small, align=left, anchor=west},
}
\begin{document}
\sffamily\hyphenpenalty=10000\exhyphenpenalty=10000

\begin{tikzpicture}
\tikzset{tc/.style={box, minimum width=31mm, minimum height=50mm, text width=28mm, font=\footnotesize, align=left, anchor=north}}
\foreach \i/\x in {1/0,2/36,3/72,4/108,5/144} \node[tc] (t\i) at (\x mm,0) {};
\node[tc, fill=band] at (t1.north) {\textbf{উপকরণ}\\[1mm]মূল দল ও ফেলো\\উন্মুক্ত নথি ও ডেটা\\যৌথ কম্পিউটিং\\অর্থায়ন ও সমঝোতা স্মারক};
\node[tc] at (t2.north) {\textbf{কার্যক্রম}\\[1mm]কর্পাস ও অনুসন্ধান নির্মাণ\\বেঞ্চমার্ক তৈরি\\গেটের অধীনে পাইলট\\কর্মকর্তা প্রশিক্ষণ};
\node[tc] at (t3.north) {\textbf{উৎপাদ}\\[1mm]P1 অনুসন্ধান ও P2 ডেটা\\BanglaPolicyQA\\সনদ, ERRC, নিবন্ধন\\প্রশিক্ষিত কর্মকর্তা};
\node[tc, fill=accentlight, draw=accent, line width=0.9pt] at (t4.north) {\textbf{ফলাফল}\\[1mm]দ্রুত, উদ্ধৃতিযুক্ত প্রমাণ\\নিয়ন্ত্রিত, আস্থাভাজন এআই\\জাতীয় বাংলা এআই সক্ষমতা};
\node[tc, fill=oklight, draw=ok] at (t5.north) {\textbf{প্রভাব}\\[1mm]প্রমাণে দৃঢ় জননীতি\\সরকারি খাতের এআই-তে জন-আস্থা};
\foreach \i/\j in {1/2,2/3,3/4,4/5} \draw[arr, line width=1pt] (t\i.east |- 0,-25mm) -- (t\j.west |- 0,-25mm);
\node[soft, minimum width=175mm, text width=170mm, font=\footnotesize, align=left] (as) at (72mm,-72mm) {\textbf{মূল অনুমান:} সংস্থাগুলো কর্মীদের সময় দেয়; ডেটা প্রবেশাধিকার মেলে; ERRC স্বাধীন থাকে; প্রশিক্ষিত কর্মীরা পদে থাকেন; সময়মতো অর্থছাড় হয়।};
\draw[darr] (t5.south) -- ++(0,-8mm) -| node[lbl, pos=0.25]{মূল্যায়ন, শিক্ষা ও ব্লুপ্রিন্ট সংশোধন} (t1.south);
\end{tikzpicture}
\end{document}
